% docs/resumen.tex
\section{Resumen}

En el departamento del Meta (Colombia), las condiciones ambientales aceleran la degradación del concreto y elevan los costos de mantenimiento. Ante la ausencia de herramientas locales para apoyar la dosificación y selección de aditivos, este proyecto desarrolla DIAPCE: una aplicación multiplataforma en Flutter \textbf{dirigida exclusivamente a los estudiantes del semillero de investigación y a los participantes en prácticas de laboratorio de materiales de la Universidad Cooperativa de Colombia (UCC), sede Meta}. DIAPCE no es una herramienta de uso industrial ni general; su propósito es formativo e institucional. A partir de un conjunto de datos regional y entradas como resistencia objetivo, temperatura, humedad y condiciones de exposición, el sistema sugiere una dosificación base y clasifica el uso de aditivos (\textbf{P0}: sin~aditivo; \textbf{P1}--\textbf{P3}: impermeabilizante al 2, 3 y~4\,\%; \textbf{PP1}--\textbf{PP3}: plastificante al 0{,}2, 0{,}4 y~0{,}6\,\%), con proyecciones a 7, 14 y 28 días, en concordancia con ACI 211.1 y normas relacionadas. La solución funciona sin conexión mediante SQLite y prioriza la trazabilidad de los ensayos y ejercicios académicos (registro de insumos, resultados y versión del modelo), facilitando prácticas reproducibles y comparables en el aula y el laboratorio. DIAPCE busca fortalecer competencias en diseño de mezcla, análisis de resultados y cultura de registro en los estudiantes de la UCC, y potencialmente contribuir a un uso más eficiente de materiales en el contexto de sus prácticas formativas.

\vspace{0.5em}
\noindent\textbf{Palabras clave:} concreto estructural, dosificación, aditivos (\textbf{P0}–\textbf{PP3}: impermeabilizante 0--4\,\% y plastificante 0{,}2--0{,}6\,\%), ACI 211.1, Flutter, SQLite, Meta (Colombia), 7/14/28 días, laboratorios, semilleros de investigación, docencia.
